% Folder 84, batch 10 (archivists' pages 181-192).
% Pass: Opus 5.5 (claude-opus-5-5), 2026-10-09 — first pass, unchecked against the pages by a human.
% Revised: Opus 5.5 (claude-opus-5-5), 2026-10-10 — re-read on the 700 dpi tiles; 6 \ill and 2 \uncertain resolved, 28 remain (18 \ill, 10 \uncertain); one misreading (« calculable », p. 189) reverted to \ill.
%
% Even pages 182, 184, 186, 188, 190, 192 are versos of a typescript of
% Grothendieck's own non-mathematical writing (typed pages 495-500, yin and
% yang); not the mathematics in hand, so skipped without placeholder.
\documentclass[11pt,a4paper]{article}
\input{../preamble/grothendieck}

\folder{84}
\batch{10}
\pages{181}{192}
\foldertitle{[Formes quadratiques] : notes manuscrites (s.d.).}
\dating{[à partir de 1982-vers 1986]}
\watermark{Édition de démonstration}

\begin{document}

\note{Ces six feuillets sont des brouillons de calcul sans titre ni texte suivi ; le contexte (algèbres $A_1$, $A_2$, modules $E_1$, $E_2$ munis de droites $L_1$, $L_2$, accouplement $A\times E\to\underline{O}$, produit $*$) vient d'avant le lot, et le calcul s'arrête sur p.~191 sans conclusion visible. Notations fixées pour le lot : $\underline{O}$ pour le symbole souligné qu'il écrit pour l'anneau de base, $k$ pour le corps (ou anneau) des scalaires.}

\page{181}
\[
(f+\sigma(f))\,\varphi(x) \;=\; \langle f+\sigma(f),\, x\rangle
\]
\[
\varphi(f)\,(x+\sigma(x)) \;=\; \langle f,\, x+\sigma x\rangle
\]
\[
\langle (1+\sigma)f,\, x\rangle
\]

\struck{$\tfrac12\,\mathrm{Tr}$} \ill{} \uncertain{d'après l'exemple}

\[
\text{\struck{$A_1$}}\otimes A_2)\otimes(E_1\otimes E_2)\longrightarrow \underline{O}
\]
\[
\mapsto\;\bigl[(A_1\otimes_0 A_2)\otimes(E_1\otimes E_2)\bigr]\oplus
(A_1\otimes_0 A_2)\otimes(E_1\otimes_0 E_2)
\]
\marginal{$[(A_1\otimes_0 A_2)\otimes(E_1\otimes_0 E_2)$}

Tableau à deux colonnes, indices $1$ et $2$ :
\[
\begin{array}{c|c}
f_1 & f_2\\
\varphi_1(f_1)\in L_1^{\vee} & \varphi_2(f_2)\\
\hline
x_1,\ T_1 & x_2,\ T_2\\
\varphi_1(x_1)\in k & \varphi_2(x_2)\\
\pi_1(x_1)\in L_1 & \pi_2(x_2)
\end{array}
\]
\note{dans la colonne $1$, la case du bas est encadrée et marquée $E_1$ à gauche, avec des flèches de $x_1$ vers $\varphi_1(x_1)$ et vers $\pi_1(x_1)$ ; au-dessus de la ligne, la case $\varphi_1(f_1)$ est surchargée et en partie biffée.}

\[
\langle (1+\sigma)f,\, x\rangle =
\]
\[
\tfrac12\bigl[(f+\sigma f)\bigr]\bigl[
\]
\[
\langle f+\sigma f,\, x\rangle
\]

\page{183}
\struck{$\lambda=\varphi(x)$}

\[
A\simeq (A_1\otimes A_2)\oplus \cdots\;\big/\;(A_1\oplus_k A_2)
\]
pour
\[
A_1\oplus_k A_2 \xrightarrow{\ \mathrm{Tr}_1,\ \mathrm{Tr}_2\ } k,
\qquad
A_1\oplus_k A_2 \xhookrightarrow{\ \text{can.}\ } A_1\otimes A_2
\]
\note{les deux flèches partent de $A_1\oplus_k A_2$ en diagonale ; tout ce premier bloc est barré au crayon d'une grande croix.}

\[
E\simeq (E_1\otimes E_2)\oplus(L_1\otimes L_2)\;\big/\;
(E_1\otimes L_2)+_{L_1\otimes L_2}(L_1\otimes E_2)
\]

$A\times E\longrightarrow \underline{O}$ \quad couplage,

\[
0\to k\to A\to L_1^{\vee}\otimes L_2^{\vee}\to 0
\]
\[
0\to L_1\otimes L_2\to E\to \underline{O}\to 0
\]

\[
\langle f_1\otimes f_2,\ \cdots\rangle
\;=\;
\text{\struck{$\langle\varphi_1(x_1),u_1\rangle\langle\varphi_2(x_2),u_2\rangle$}}
\;\varphi_1(f_1)\varphi_2(f_2)\cdot u
\]
avec $f_1\otimes f_2\in A_1\otimes A_2$, $u\in L_1\otimes L_2$ ; sous la ligne : $\varphi(x_1\otimes x_2)$.

\[
\langle\lambda,\ x_1\otimes x_2\rangle = \lambda\,\psi_1(f_1)\psi_2(f_2)
= \lambda\,\psi(x_1\otimes x_2)
\]
\[
\langle\lambda,\ u\rangle = 0
\]
\[
\langle f_1\otimes f_2,\ x_1\otimes x_2\rangle
\]
\note{sous cette dernière ligne, \struck{$f_1\otimes 1$}, biffé.}

\struck{expr. bil. en $f_1\otimes f_2$, $x_1\rangle$ $\langle f_2, x_2\rangle$} \struck{\uncertain{nul sur}} \struck{\ill{}}, $E$)
(\struck{$A_1\otimes A_2$} $A_1\oplus A_2$, \uncertain{et non}

si $f\in A_1$ ou $A_2$, \ill{} \ill{} \ill{} \ill{}

si $x\in E_1\otimes L_2+L_1\otimes E_2$, \ill{} \ill{}
\note{sous ces deux dernières lignes, des annotations serrées : $\mathrm{T}(f)\,\psi(x)$ et $\varphi(f)\cdot$\ill{}, illisibles en partie.}

\page{185}
\[
\begin{array}{cccc}
1, & u_1, & u_2, & u_1u_2\\
\downarrow & \downarrow & \downarrow & \downarrow\\
2 & \cdot, & \cdot & 2u_1u_2=v
\end{array}
\qquad \sigma_1\sigma_2
\]
\note{les deux images du milieu sont cerclées et surchargées ; on devine $2u_1$, $2u_2$ sous la surcharge.}

\[
v\longmapsto 2u_1u_2
\]

\[
-u_1^2=c_1\cdot 1,\qquad -u_2^2=c_2\cdot 1,\qquad u_1^2u_2^2=c_1c_2
\]

\[
x_1\otimes x_2\longmapsto x_1\otimes x_2+\overline{x}_1\otimes\overline{x}_2
\]

\begin{tikzcd}[column sep=small, row sep=large, nodes={font=\small}]
E_1\otimes E_2 \arrow[r] \arrow[rr, bend left=20] & E_1\wedge E_2 \arrow[r, Rightarrow, "?"] & E_1\otimes E_2 \arrow[d, leftrightarrow, "\text{can}"] \\
E_1\otimes_0 E_2 \arrow[u, "\text{can}"] \arrow[r, "\pi"'] & L_1\otimes L_2 \arrow[u] \arrow[ur] \arrow[r, "{\mathrm{id}}"'] & L_1\otimes L_2
\end{tikzcd}

\note{sur la flèche $L_1\otimes L_2\to E_1\otimes E_2$ (diagonale) : $u_1\otimes u_2\mapsto u_1\otimes u_2$, puis « i.e. \ill{} » ; la flèche horizontale du bas est marquée $\mathrm{id}$.}

\[
E_1\otimes L_2 \ni x_1\otimes u_2
\]
\[
\pi_1(x_1)\otimes u_2=(2x_1-\overset{\psi_1(x_1)}{T_1})\otimes u_2
\]
\[
x_1\otimes u_2+(\psi_1(x_1)T_1-x_1)\otimes(-u_2)
\]
\[
=2x_1\otimes u_2-\psi_1(x_1)T_1\otimes u_2
\]
\marginal{ok}

\[
(x_1+u_1)*(x_2+u_2)=(\sigma+\overline{\sigma})\bigl[(x_1+u_1)\otimes(x_2+u_2)\bigr]
\]
\[
=x_1*x_2+\pi_1(x_1)\otimes u_2+u_1\otimes\pi_2(x_2)+\underbrace{u_1*u_2}_{=\,2u_1\otimes u_2}
\]
\[
b_1\otimes u_2+u_2\otimes b_1+2u_1\otimes u_2
\]

\struck{$b_1u_2+b_2u_1+u_1$}

\[
(x_1+u_1)*(x_2+u_2) = x_1*x_2+\underbrace{(b_1u_2+b_2u_1+2u_1u_2)}\ \in E_1\wedge E_2
\]
\note{sous $x_1*x_2$ une flèche descend vers $x_1*x_2 = \mathrm{Tr}_{E_1\otimes E_2/E_1\wedge E_2}(x_1\otimes x_2)$ ; sous $E_1\wedge E_2$ une flèche descend vers $E_1\otimes E_2$.}

\page{187}
Considérons $\langle u_1*u_2,\ x_1*x_2\rangle=$

\struck{$x_1*x_2$}
\[
A_1\otimes A_2\to (A_1\otimes A_2)\oplus(A_1\oplus_0 A_2)\to A
\]
\[
A_1\otimes A_2\xrightarrow{\ \text{can}\ }(A_1\otimes A_2)\oplus(A_1\oplus_0 A_2)
\xrightarrow{\ \text{somme}\ } A
\]
\note{le second $\oplus_0$ de cette ligne porte un $0$ souligné.}

\struck{$E_1\otimes E_2$} \qquad $\longrightarrow E$

\[
\text{\struck{$A^{\vee}\simeq A_1\otimes A_2$}}\times\cdots\;\big/\;\mathrm{Im}(A_1\oplus_k A_2)
\]
\struck{$E^{\vee}\simeq$ Formes bilin. sur $A_1\times A_2$ telles que}

\[
A^{\vee}\simeq(f,\overset{\lambda}{\cdots})\quad\ldots
\]
\begin{itemize}
  \item $f$ forme bil. sur $A_1\times A_2\to k$, \quad $\lambda\in k$
  \item \struck{$g_1 : A_1\to k$}
  \item \struck{$g_2 : A_2\to k$}
\end{itemize}

\uncertain{définition} \struck{\ill{}} \quad
\struck{a) $g_1(1)=g_2(1)$ i.e. $g$ provient d'une forme sur $A_1+A_2\subset A_1\otimes A_2\to k$}

\struck{b) $f\,|\,(A_1+A_2)=2g$} i.e.

\[
f(x_1,1)=2g_1\qquad x_1\in A_1
\]
\[
f(1,x_2)=2g_2\qquad x_2\in A_2
\]
\note{ces deux lignes sont barrées d'un trait en boucle.}

\[
f(x_1)=\lambda\,\mathrm{Tr}_1(x_1)
\]
\[
f(x_2)=\lambda\,\mathrm{Tr}_2(x_2)
\]
\[
0\to L_1\otimes L_2\to E\xrightarrow{\ \psi\ }\underline{O}\to 0
\]
avec $\check{a}\in E$ sous $E$.

\struck{$E\to\check{A}$} \qquad \struck{$\lambda=\varphi(x)$} \qquad $E$

\page{189}
\struck{\ill{}} $u$ \ill{} en termes de $u_1$, $u_2$ ?

\struck{$u=u_1$ \uncertain{Comparons}}

\uncertain{Comment} \struck{\uncertain{Considérons}} \quad On \uncertain{doit avoir}
\[
u_1\overset{\xi}{*}u_2=\underset{\sim}{\lambda}\,u+\underset{\sim}{\mu}
\qquad\text{trouver } \lambda,\mu
\]
\[
\mathrm{Tr}_{A/k}(u_1*u_2)=-\lambda\beta+2\mu
\]
\[
N_{A/k}(u_1*u_2)=\underline{\lambda^2}c+\underline{\mu^2}-\underline{\lambda\mu\beta}
\]

\begin{itemize}
  \item base de $A$ : \quad $1,\ u$
  \item base de $E$ : \quad \struck{$e=e_1\otimes e_2$, $x=x_1*x_2$} \quad $x=x_1*x_2$, $e=e_1\otimes e_2$
  \item (\uncertain{duale})
\end{itemize}
\marginal{\struck{$\mathrm{Tr}_{A_1\otimes A_2}$}}

\[
\langle 1,x\rangle=1\qquad \langle 1,e\rangle=0
\]
\[
\langle u,e\rangle=1\qquad \langle u,x\rangle=0
\]

\[
\lambda=\text{\struck{$\langle u,e\rangle$}}\;\langle u_1*u_2,\ \overbrace{e_1\otimes e_2}^{e}\rangle
\]
\[
\mu=\langle u_1*u_2,\ \underbrace{x_1*x_2}_{x}\rangle
\]

$u_1*u_2$

\begin{tikzcd}[column sep=large, row sep=small, nodes={font=\small}]
A_1\otimes A_2 \arrow[r, "*"] \arrow[rr, bend left=20] & A \arrow[r, "\xi"] & A_1\otimes A_2 \\
E_1\otimes E_2 \arrow[d] & E \arrow[l, "\xi"'] \arrow[d] & E_1\otimes E_2 \arrow[l, "(*)"'] \arrow[d] \\
\underline{O} & \underline{O} & \underline{O}
\end{tikzcd}

\note{les colonnes sont appariées par de petites croix « $\times$ » entre les deux lignes (dualité $A\times E$), et un « $\wedge$ » est posé entre $A_1\otimes A_2$ et le $E_1\otimes E_2$ de droite ; l'étiquette de $A\to A_1\otimes A_2$ est cerclée, avec « $(\mathrm{Tr})$ » au-dessus ; au-dessus de la flèche courbe, \struck{$x_1\otimes x_2$} biffé et $\mathbb{F}_2\to\mathbb{F}_2\times\mathbb{F}_2$.}

\[
\text{\struck{$(x,y)\mapsto x\otimes\cdots$}}\qquad
{}^t(\mathrm{id}+\sigma_1\otimes\sigma_2)=\mathrm{id}+({}^t\sigma_1)\otimes({}^t\sigma_2)
\]
avec ${}^t\sigma_1=\sigma_1$, ${}^t\sigma_2=\sigma_2$.

$u\in A$ défini par
\[
\begin{cases}
\langle u,\ e_1\otimes e_1\rangle=1\\
\langle u,\ x_1*x_2\rangle=0
\end{cases}
\]
\note{« $e_1\otimes e_1$ » est sur la page ; on attendrait $e_1\otimes e_2$.}

\page{191}
\[
u_1\qquad u_1^2+b_1u_1+c_1=0
\]
\[
u_2\qquad u_2^2+b_2u_2+c_2=0
\]
\[
u_1*u_2=u_1\otimes u_2+(u_1+b_1)\otimes(u_2+b_2)=2u_1u_2+b_2u_1+b_1u_2+\underline{\underline{b_1b_2}}
\]
\struck{\ill{} \ill{} \ill{}} \ill{} \uncertain{bien des}

$A$ : avec $1$ \qquad \ill{} si $2=0$, $b_1=b_2=0$

\[
A_1\simeq\underline{O}\oplus L_1^{\vee}\qquad A_2=\underline{O}\oplus L_2^{\vee}
\]
\[
E_1\simeq\underline{O}+L_1,\qquad E_2\simeq\underline{O}+L_2
\]
\[
E_1\otimes E_2\to E_1\wedge E_2\to E_1\otimes E_2
\]
\[
x_1\otimes x_2\qquad x_1*x_2\qquad \mathrm{Tr}(x_1\otimes x_2)
\]

\struck{$A\simeq L_1^{\vee}\otimes L_2^{\vee}$}
\[
A\simeq\underline{O}\oplus L_1^{\vee}\otimes L_2^{\vee}
\]

$e_1\in L_1$ base de $L_1$, \qquad $e_2\in L_2$ base de $L_2$

\[
b_1=\beta_1e_1\qquad c_1=\gamma_1e_1^{\otimes 2}
\]
\[
b_2=\beta_2e_2\qquad c_2=\gamma_2e_2^{\otimes 2}
\]

\[
A\simeq\underline{O}+e_1^{*}\ldots\qquad u_1=e_1^{*}\qquad u_2=e_2^{*}
\]
\note{après $\underline{O}+$, un symbole surchargé, lu $e_1^{*}$ sous réserve.}

\struck{$u_1^2+\beta_1u_1 = e_1^{*}$}
\[
u_1^2+\beta_1u_1+\gamma_1=0
\]
\[
u_2^2+\beta_2u_2+\gamma_2=0
\]
\note{sous la seconde équation, une ligne barrée et surchargée, où l'on devine $u_1^2+(\beta_1\ldots)\,u_1$ ; le $u_2^2$ de la seconde équation est écrit sur un $u_1$.}

$u=2uu'+$\note{le $2$ est surchargé ; rien n'est écrit après le $+$.}

\[
x_1,\ e_1\ \longleftrightarrow\ u_1\in A_1\qquad u_1^2+\beta_1u_1+\gamma_1=0
\]
\[
\begin{cases}
e_1=\varphi(u_1)\\
\langle u_1,\ x_1\rangle=0\\
\varphi(x_1)=1
\end{cases}
\]
\note{la notation $\langle u_1, x_1\rangle$ est surchargée ; à gauche de l'accolade, un mot biffé.}
\[
(x_2,\ e_2)\ \longleftarrow\ u_2\in A_2\qquad u_2^2+\beta_2u_2+\gamma_2=0
\]
\[
(x_1*x_2,\ e_1\otimes e_2)\ \longleftrightarrow\ u\in A.
\]
\[
u^2+\beta u+\gamma=0
\]

\end{document}
